\SetPractica{Determinación de ácido cítrico}{Determinación del porcentaje de ácido cítrico y ácido acético en jugos de frutas y vinagre}

\section{Introducción teórica}



\section{Objetivos}

\begin{itemize}
    \item Preparar soluciones químicas adecuadas para utilizar en un método analítico volumétrico
    \item Observar una reacción de neutralización ácido-base y establecer su balance estequiométrico
    \item Determinar el porcentaje de ácido cítrico y de ácido acético en sustancias utilizadas cotidianamente en el hogar
\end{itemize}

\section{Materiales, equipos y reactivos}

\begin{table}[H]
\centering{\small\setstretch{1.25}
\begin{tabular}{|m{0.3\linewidth}|m{0.3\linewidth}|m{0.3\linewidth}|} \hline
    
    \thead{Equipos} & \thead{Materiales} & \thead{Reactivos} \\ \hline
    
    {Potenciómetro o tiras de pH} 
    
    &
    
    {Cristalería para preparación de reactivos}
    {Soporte universal \par}
    {Bureta \par}
    {Vasos de precipitado \par}
    {Pipetas}
    
    & 
    
    {Solución de hidróxido de sodio (\ce{NaOH}) al 0.1N \par}
    {Fenolftaleína al 0.1\% \par}
    {Jugo de naranja natural \par}
    {Jugo de naranja procesado \par}
    {Jugo de limón natural \par}
    {Vinagre comercial}
    
    \\ \hline
\end{tabular}}\end{table}

\section{Procedimientos}

\begin{enumerate}
    \item Agitar vigorosamente las muestras antes de usarlas.
    \item Medir el pH de las muestras con el potenciómetro o tiras de pH
    \item Llenar la bureta con la solución de hidróxido de sodio 0.1N.
    \item Transferir con una pipeta 1ml de cada una de las muestras (alícuota) a un vaso de precipitado de 50ml
    \item A cada vaso con muestra añadir (con una pipeta) 3ml de agua destilada y agregar una a dos gotas de fenolftaleína al 0.1\% como indicador.
    \item Colocar el vaso de precipitado bajo la bureta dejando caer gota a gota el \ce{NaOH} agitando continuamente el vaso hasta que la solución adquiera un color rosado; anotar el volumen de \ce{NaOH} gastado. Para cada muestra la determinación se debe hacer por triplicado. 
    \item Repetir el procedimiento con cada una de las muestras, anotando los volúmenes de la solución de \ce{NaOH} utilizados en cada titulación.
\end{enumerate}